% Source: adjoint_operator.tex
\begin{notechapter}
\chapter{The Adjoint Operator and the Dual Space}\label{ch:adjoint-operator}

\section{Introduction}
This document clarifies the concepts of the adjoint operator and the dual operator, addressing the subtle differences and connections between these fundamental ideas in linear algebra and functional analysis. We will emphasize the role of the inner product and the Riesz Representation Theorem in linking purely algebraic notions (dual spaces) to geometric ones (adjoints).

\section{The Foundation: Transpose and Inner Products}

\begin{definition}[Transpose]
For a real matrix $A \in \mathbb{R}^{n \times m}$, representing a map $A: \mathbb{R}^m \to \mathbb{R}^n$, the transpose $A^T \in \mathbb{R}^{m \times n}$ is defined by swapping rows and columns: $(A^T)_{ji} = A_{ij}$.
\end{definition}

The transpose has a key property related to the standard dot product in Euclidean space.

\begin{theorem}[Fundamental Property of Transpose]
For any real matrix $A \in \mathbb{R}^{n \times m}$ and vectors $x \in \mathbb{R}^m$, $y \in \mathbb{R}^n$:
\[
(Ax) \cdot y = x \cdot (A^T y)
\]
where $\cdot$ denotes the standard dot product.
\end{theorem}

This property highlights that the transpose is intrinsically linked to the geometric structure provided by the inner product.

\begin{definition}[Inner Product Space]
An inner product space is a vector space $V$ over $\mathbb{F}$ (where $\mathbb{F} = \mathbb{R}$ or $\mathbb{C}$) equipped with an inner product $\langle \cdot, \cdot \rangle$ satisfying:
\begin{enumerate}
    \item Conjugate symmetry: $\langle x, y \rangle = \overline{\langle y, x \rangle}$
    \item Linearity in the first argument: $\langle ax + by, z \rangle = a\langle x, z \rangle + b\langle y, z \rangle$
    \item Positive-definiteness: $\langle x, x \rangle \geq 0$, and $\langle x, x \rangle = 0 \iff x = 0$
\end{enumerate}
\end{definition}

\section{The Dual Operator: A Purely Algebraic Concept}

The dual operator exists for any linear map between vector spaces, requiring no inner product structure.

\begin{definition}[Dual Space]
For a vector space $V$ over $\mathbb{F}$, the dual space $V^*$ is the set of all linear functionals:
\[
V^* = \{ \phi : V \to \mathbb{F} \mid \phi \text{ is linear} \}.
\]
\end{definition}

\begin{definition}[Dual Operator]
For a linear operator $T : V \to W$ between vector spaces $V$ and $W$, the dual operator $T^* : W^* \to V^*$ is defined by:
\[
T^*(\phi)(v) = \phi(T(v)) \quad \text{for all } \phi \in W^*, v \in V.
\]
Equivalently, $T^*(\phi) = \phi \circ T$. The dual operator maps functionals to functionals.
\end{definition}

In finite dimensions with fixed bases, the matrix representation of the dual operator $T^*$ (with respect to the dual bases) is exactly the transpose $A^T$ of the matrix $A$ representing $T$ --- also in the complex case: the dual operator is linear and involves no complex conjugation. The conjugate transpose $A^\dagger$ appears only for the \emph{adjoint}, which uses an inner product.

\section{The Adjoint Operator: A Geometric Realization}

The adjoint operator brings in the geometric structure of an inner product.

\begin{definition}[Adjoint Operator]
Let $V$ and $W$ be Hilbert spaces (for instance finite-dimensional inner product spaces) over the same field $\mathbb{F}$. For a bounded linear operator $A : V \to W$, the adjoint operator $A^* : W \to V$ is the unique operator satisfying:
\[
\langle Av, w \rangle_W = \langle v, A^*w \rangle_V \quad \text{for all } v \in V, w \in W.
\]
The adjoint maps vectors to vectors, unlike the dual operator which maps functionals to functionals.
\end{definition}

\begin{example}[Derivative Operator]
Consider the derivative operator $D = \frac{d}{dx}$ on $L^2([a, b])$ with functions vanishing at boundaries. Using integration by parts for the inner product $\langle f, g \rangle = \int_a^b f(x)g(x)dx$:
\[
\langle Df, g \rangle = \int_a^b f'(x)g(x)dx = [f(x)g(x)]_a^b - \int_a^b f(x)g'(x)dx
\]
The boundary terms vanish, leading to $\langle Df, g \rangle = \langle f, -Dg \rangle$, so $D^* = -D$.

Strictly speaking, $D$ is \emph{not} a bounded operator on $L^2([a,b])$ (e.g.\ $\|\sin(kx)\|$ stays bounded while $\|D\sin(kx)\|$ grows like $k$), so the definition above does not apply directly: $D$ is only defined on a dense subspace of differentiable functions, and the computation shows that $-D$ is its \emph{formal} adjoint. Making this rigorous requires the theory of unbounded operators, in which the domain of the adjoint is part of its definition. Note also that it suffices that $f$ (or $g$) vanish at the boundary for the boundary term to disappear.
\end{example}

\section{The Riesz Theorem: Connecting the Concepts}

The Riesz Representation Theorem is the bridge between the algebraic dual space and the geometric adjoint operator.

\begin{theorem}[Riesz Representation Theorem]
For a Hilbert space $V$, the map $R_V : V \to V^*$ defined by $R_V(w) = \langle \cdot, w \rangle$ is a conjugate-linear isomorphism (or a linear isomorphism if defined as $R_V(w) = \langle w, \cdot \rangle$ in real spaces). This means every linear functional $\phi \in V^*$ can be uniquely represented as $\phi(v) = \langle v, w \rangle$ for some $w \in V$.
\end{theorem}

The adjoint operator is the "concrete" realization of the dual operator using these isomorphisms. The two concepts are related by the following commutative diagram, where $A^*_{\text{adjoint}}$ makes the path from $W$ to $V$ below equivalent to the path through the dual spaces:

\subsection{General Commutative Diagram for Hilbert spaces $V$ and $W$}
\[
\begin{tikzcd}[row sep=large, column sep=large]
V \arrow[d, "R_V"'] \arrow[r, "A"] & W  \\
V^* \arrow[d, "R_V^{-1}"] & W^* \arrow[l, "{A^*_{\text{dual}}}"'] \arrow[u, "R_W^{-1}"'] \\
V & W \arrow[l, "A^*_{\text{adjoint}}"] \arrow[u, "R_W"']
\end{tikzcd}
\]

The relationship is given by the composition:
\[
A_{\text{adjoint}}^{*} = R_{V}^{-1} \circ A_{\text{dual}}^{*} \circ R_{W}
\]

\subsection{The Finite-Dimensional Case ($\mathbb{R}^m \to \mathbb{R}^n$)}

When considering a map $A: \mathbb{R}^m \to \mathbb{R}^n$ with standard bases and dot products, we have $V = \mathbb{R}^m$, $W = \mathbb{R}^n$. The dual spaces are naturally isomorphic to the original spaces ($\mathbb{R}^m \cong (\mathbb{R}^m)^*$, $\mathbb{R}^n \cong (\mathbb{R}^n)^*$) via the Riesz map which becomes the identity operation.

\begin{lemma}
Let $V = \mathbb{R}^m$ with the standard basis $\{e_1, \dots, e_m\}$ and the standard dot product. The Riesz isomorphism $R_V : V \to V^*$, defined by $R_V(w) = \langle \cdot, w \rangle$, is the identity map when we identify $V$ with $V^*$ using the standard dual basis.
\end{lemma}

\begin{proof}
The standard dual basis $\{\varepsilon_1, \dots, \varepsilon_m\}$ for $V^*$
is defined by $\varepsilon_i(e_j) = \delta_{ij}$.\\
For any vector $w = \sum_{i=1}^m w_i e_i \in V$, the functional $R_V(w)$ acts on a vector $v = \sum_{j=1}^m v_j e_j$ as:
\[
R_V(w)(v) = \langle v, w \rangle = \sum_{j=1}^m v_j w_j.
\]
Now, consider the functional $\phi_w = \sum_{i=1}^m w_i \varepsilon_i \in V^*$. Its action on $v$ is:
\[
\phi_w(v) = \sum_{i=1}^m w_i \varepsilon_i\left(\sum_{j=1}^m v_j e_j\right) = \sum_{i=1}^m w_i v_i = \sum_{j=1}^m v_j w_j.
\]
Since $R_V(w)(v) = \phi_w(v)$ for all $v \in V$, we have $R_V(w) = \phi_w$. Under the natural identification $e_i \leftrightarrow \varepsilon_i$, the vector $w$ corresponds exactly to the functional $\phi_w$. Therefore, $R_V$ is the identity map.
\end{proof}

In this specific, highly practical scenario, the commutative diagram simplifies conceptually as follows:

\[
\begin{tikzcd}[row sep=large, column sep=large]
\mathbb{R}^m \arrow[d, equal, "R_V \approx I"] \arrow[r, "A" {above}, shift left=0.5ex] & 
\mathbb{R}^n \arrow[d, equal, "R_W \approx I"] \arrow[l, "A^*_{\text{adjoint}}" {below}, shift left=0.5ex]\\
(\mathbb{R}^m)^* & (\mathbb{R}^n)^* \arrow[l, "{A^*_{\text{dual}} = A^T}" {below}]
\end{tikzcd}
\]

\begin{remark}
In $\mathbb{R}^m \to \mathbb{R}^n$ with the standard basis and 
dot product, the Riesz isomorphism is represented by the identity 
matrix. This is why we can safely abuse notation and say 
$A^*_{\text{adjoint}} = A^T$, as the Riesz maps are implicitly 
identity transformations in coordinate computations. 
In complex spaces $\mathbb{C}^m \to \mathbb{C}^n$, 
the matrix of the adjoint is the conjugate transpose, $A^\dagger$.
\end{remark}

\section{Summary of Relationships}

\begin{itemize}[leftmargin=*]
    \item \textbf{Dual Operator ($T^*$ or $T^*_{\text{dual}}$):} A purely algebraic
    concept for any vector spaces $V, W$. Maps $W^* \to V^*$.
    Always exists.
    \item \textbf{Adjoint Operator ($A^*$ or $A^\dagger$):} A geometric
    concept for inner product spaces $V, W$. Maps $W \to V$.
    Exists via the Riesz map relating it to the dual operator.
    \item \textbf{Transpose/Conjugate Transpose ($A^T/A^\dagger$):} 
    The matrix representation of the dual operator (relative to 
    dual bases) which becomes the matrix representation of the 
    adjoint operator when standard orthonormal bases are used.
\end{itemize}
\end{notechapter}
